\documentclass[UTF8]{ctexart}

% 支持从项目根目录或当前文件目录编译。
\makeatletter
\def\input@path{{../}}
\makeatother
\usepackage{researchnotes}

\title{三角函数公式与解三角形}
\author{}
\date{}

\begin{document}
\maketitle
\tableofcontents

\section{引言与学习路线}

三角函数把角的大小与长度、坐标联系起来。一方面，它们用于求解三角形中的边、角和面积；另一方面，角的相加、相减与倍增又产生一组相互联系的恒等式，使三角表达式能够在乘积、和差、幂与有理式之间转换。

本文先用单位圆定义三角函数，建立基本性质与和差角公式，再依次推导倍角、半角、积化和差、和差化积及万能公式，最后讨论三角方程和解三角形。学习时应把和差角公式看作推导其他公式的起点，并根据问题的结构选择变形方向，而不是把每条公式孤立地记忆。

除特别说明外，角均为实数，采用\keyterm{弧度制}（radian measure）；$k$ 表示任意整数，$\mathbb Z$ 表示整数集。所有含分母的等式都只在分母非零时使用。恒等式表示在指定定义域内对每一个自变量都成立的等式，三角方程则要求找出使等式成立的特定自变量。

\section{角、单位圆与基本性质}
\label{sec:trig-basics}

\subsection{弧度与单位圆定义}

半径为 $r>0$ 的圆上，圆心角所对的弧长为 $s$ 时，该角的弧度数为 $s/r$。逆时针旋转为正，顺时针旋转为负，并允许旋转多于一周。因此
\[
  2\pi\text{ 弧度}=360^\circ,\qquad
  1^\circ=\frac{\pi}{180}\text{ 弧度}.
\]
例如 $30^\circ=\pi/6$，$45^\circ=\pi/4$，$60^\circ=\pi/3$。同一道计算中应统一角度单位。

\begin{definition}[三角函数]
在直角坐标平面内，把点 $(1,0)$ 绕原点旋转角 $x$，所得单位圆上的点记为 $P$。定义\keyterm{余弦}（cosine）与\keyterm{正弦}（sine）为 $P$ 的横、纵坐标，即
\[
  P=(\cos x,\sin x).
\]
进一步定义\keyterm{正切}（tangent）与\keyterm{余切}（cotangent）为
\[
  \tan x=\frac{\sin x}{\cos x}\quad(\cos x\ne0),\qquad
  \cot x=\frac{\cos x}{\sin x}\quad(\sin x\ne0).
\]
\keyterm{正割}（secant）与\keyterm{余割}（cosecant）分别为
\[
  \sec x=\frac1{\cos x}\quad(\cos x\ne0),\qquad
  \csc x=\frac1{\sin x}\quad(\sin x\ne0).
\]
\end{definition}

在直角三角形中，对锐角 $x$，上述定义等价于
\[
  \sin x=\frac{\text{对边}}{\text{斜边}},\qquad
  \cos x=\frac{\text{邻边}}{\text{斜边}},\qquad
  \tan x=\frac{\text{对边}}{\text{邻边}}.
\]
这是因为把斜边缩放为 $1$ 后，两条直角边恰好成为单位圆上相应点的坐标。单位圆定义进一步把这些比值推广到了锐角之外。

\subsection{定义域、符号、周期与奇偶性}

正弦与余弦的定义域均为 $\mathbb R$，值域均为 $[-1,1]$。正切的定义域为 $x\ne\pi/2+k\pi$，余切的定义域为 $x\ne k\pi$，二者值域均为 $\mathbb R$。

按第一、第二、第三、第四象限的顺序，$\sin x$ 的符号依次为 $+,+,-,-$，$\cos x$ 的符号依次为 $+,-,-,+$，$\tan x$ 的符号依次为 $+,-,+,-$。终边落在坐标轴上的角应直接按坐标计算，不归入象限内的符号判断。

绕圆一周后坐标不变，关于横轴反射使纵坐标变号，因此
\begin{align}
  \sin(x+2k\pi)&=\sin x,&\cos(x+2k\pi)&=\cos x,
  \label{eq:trig-period}\\
  \sin(-x)&=-\sin x,&\cos(-x)&=\cos x.
  \label{eq:trig-parity}
\end{align}
正弦是奇函数，余弦是偶函数；它们的最小正周期均为 $2\pi$。正切与余切都是奇函数，最小正周期均为 $\pi$，周期等式在各自定义域内成立。

\subsection{基本恒等式与特殊角}

单位圆方程给出\keyterm{平方关系}（Pythagorean identity）
\begin{equation}\label{eq:trig-pythagorean}
  \boxed{\sin^2x+\cos^2x=1}.
\end{equation}
这里 $\sin^2x$ 表示 $(\sin x)^2$，不表示 $\sin(x^2)$。分别除以 $\cos^2x$ 和 $\sin^2x$ 得
\begin{equation}\label{eq:trig-sec-csc}
  1+\tan^2x=\sec^2x,\qquad
  1+\cot^2x=\csc^2x,
\end{equation}
第一式要求 $\cos x\ne0$，第二式要求 $\sin x\ne0$。此外，$\tan x\cot x=1$ 要求 $\sin x\cos x\ne0$。

常用特殊角的值如下。$\pi/4$ 的值可由等腰直角三角形得到，$\pi/6$ 与 $\pi/3$ 的值可由等边三角形作高得到；坐标轴上的值直接由单位圆定义得到。
\begin{center}
\begin{tabular}{c c c c c c}
\toprule
$x$ & $0$ & $\pi/6$ & $\pi/4$ & $\pi/3$ & $\pi/2$ \\
\midrule
$\sin x$ & $0$ & $1/2$ & $\sqrt2/2$ & $\sqrt3/2$ & $1$ \\
$\cos x$ & $1$ & $\sqrt3/2$ & $\sqrt2/2$ & $1/2$ & $0$ \\
$\tan x$ & $0$ & $\sqrt3/3$ & $1$ & $\sqrt3$ & 未定义 \\
\bottomrule
\end{tabular}
\end{center}

\subsection{诱导公式与函数图像的基本信息}

单位圆上的反射与旋转给出\keyterm{诱导公式}（reduction formulas）
\begin{align}
  \sin(\pi-x)&=\sin x,&\cos(\pi-x)&=-\cos x,\label{eq:trig-supplement}\\
  \sin(\pi+x)&=-\sin x,&\cos(\pi+x)&=-\cos x,\notag\\
  \sin\left(\frac\pi2-x\right)&=\cos x,&
  \cos\left(\frac\pi2-x\right)&=\sin x,\notag\\
  \sin\left(\frac\pi2+x\right)&=\cos x,&
  \cos\left(\frac\pi2+x\right)&=-\sin x.\notag
\end{align}
例如，$\sin(7\pi/6)=-\sin(\pi/6)=-1/2$。记忆时先判断是否交换正弦、余弦，再由象限判断符号。

从单位圆坐标随角度的变化可知：$\sin x$ 在 $[-\pi/2,\pi/2]$ 上严格递增，在 $[\pi/2,3\pi/2]$ 上严格递减；$\cos x$ 在 $[0,\pi]$ 上严格递减；$\tan x$ 在 $(-\pi/2,\pi/2)$ 上严格递增。其他单调区间由周期平移得到。

对于 $y=M\sin(\omega x+\varphi)+d$，其中 $M\ne0$、$\omega\ne0$，\keyterm{振幅}（amplitude）为 $|M|$，最小正周期为 $2\pi/|\omega|$，值域为 $[d-|M|,d+|M|]$，$\varphi$ 称为\keyterm{初相位}（initial phase）。写成 $M\sin\bigl(\omega(x+\varphi/\omega)\bigr)+d$，可以识别横向平移量。

\section{和差角公式：其他公式的起点}
\label{sec:trig-addition}

\subsection{正弦、余弦的和差角公式}

\begin{theorem}[和差角公式]
对任意实数 $\alpha,\beta$，有
\begin{align}
  \sin(\alpha+\beta)&=\sin\alpha\cos\beta+\cos\alpha\sin\beta,\label{eq:trig-sin-add}\\
  \sin(\alpha-\beta)&=\sin\alpha\cos\beta-\cos\alpha\sin\beta,\label{eq:trig-sin-sub}\\
  \cos(\alpha+\beta)&=\cos\alpha\cos\beta-\sin\alpha\sin\beta,\label{eq:trig-cos-add}\\
  \cos(\alpha-\beta)&=\cos\alpha\cos\beta+\sin\alpha\sin\beta.\label{eq:trig-cos-sub}
\end{align}
这些公式称为\keyterm{和差角公式}（angle addition and subtraction formulas）。
\end{theorem}

\begin{proof}
平面旋转保持长度、垂直关系与方向。旋转角 $\alpha$ 时，两个坐标单位向量分别变成
\[
  (1,0)\longmapsto(\cos\alpha,\sin\alpha),\qquad
  (0,1)\longmapsto(-\sin\alpha,\cos\alpha).
\]
第二个像是第一个像逆时针转 $\pi/2$ 所得的垂直单位向量。旋转把向量的线性组合变为像向量的同一线性组合，故点 $(\cos\beta,\sin\beta)$ 旋转后的坐标为
\[
  (\cos\alpha\cos\beta-\sin\alpha\sin\beta,
    \sin\alpha\cos\beta+\cos\alpha\sin\beta).
\]
该点的总旋转角为 $\alpha+\beta$。比较两个坐标，得到两条和角公式；再把 $\beta$ 换为 $-\beta$，利用奇偶性即得两条差角公式。
\end{proof}

\begin{example}[求非基本特殊角的精确值]
把 $\pi/12$ 写成 $\pi/4-\pi/6$，有
\[
  \sin\frac\pi{12}
  =\frac{\sqrt2}{2}\frac{\sqrt3}{2}
   -\frac{\sqrt2}{2}\frac12
  =\frac{\sqrt6-\sqrt2}{4},\qquad
  \cos\frac\pi{12}=\frac{\sqrt6+\sqrt2}{4}.
\]
这里先把目标角拆成已知角之差，再代入公式，避免依赖小数近似。
\end{example}

\subsection{正切的和差角公式及限制}

把正弦和差角公式除以对应的余弦公式，再将分子、分母同除以 $\cos\alpha\cos\beta$，得到
\begin{equation}\label{eq:trig-tan-add}
  \boxed{\tan(\alpha\pm\beta)
  =\frac{\tan\alpha\pm\tan\beta}{1\mp\tan\alpha\tan\beta}}.
\end{equation}
这里要求 $\cos\alpha\cos\beta\ne0$，且 $1\mp\tan\alpha\tan\beta\ne0$；后一个条件在前一个条件下等价于 $\cos(\alpha\pm\beta)\ne0$。上下符号必须对应选取。

例如
\[
  \tan\frac\pi{12}
  =\frac{1-1/\sqrt3}{1+1/\sqrt3}=2-\sqrt3.
\]
但不能在 $\alpha=\beta=\pi/2$ 时使用式 \eqref{eq:trig-tan-add}，虽然 $\tan(\alpha+\beta)=\tan\pi$ 有意义，因为右侧的 $\tan\alpha$、$\tan\beta$ 没有定义。

\section{倍角、降幂、半角与三倍角公式}
\label{sec:trig-multiple}

\subsection{倍角与降幂}

在和角公式中令 $\alpha=\beta=x$，并使用式 \eqref{eq:trig-pythagorean}，得到\keyterm{倍角公式}（double-angle formulas）
\begin{align}
  \sin2x&=2\sin x\cos x,\label{eq:trig-double-sin}\\
  \cos2x&=\cos^2x-\sin^2x
  =2\cos^2x-1=1-2\sin^2x,\label{eq:trig-double-cos}\\
  \tan2x&=\frac{2\tan x}{1-\tan^2x}.\label{eq:trig-double-tan}
\end{align}
最后一式要求 $\cos x\ne0$ 且 $\tan^2x\ne1$。余弦倍角公式有三种形式：同时含正弦和余弦时用第一种，希望只保留余弦或正弦时用后两种。

反过来解出平方项，得到\keyterm{降幂公式}（power-reduction formulas）
\begin{equation}\label{eq:trig-power-reduction}
  \boxed{\sin^2x=\frac{1-\cos2x}{2},\qquad
  \cos^2x=\frac{1+\cos2x}{2}}.
\end{equation}
降幂是把同一个角的高次幂改写成倍角的低次表达式，不是把函数值近似化。

\begin{example}[四次幂的化简与范围]
对任意实数 $x$，
\begin{align*}
  \sin^4x+\cos^4x
  &=(\sin^2x+\cos^2x)^2-2\sin^2x\cos^2x\\
  &=1-\frac12\sin^22x=\frac34+\frac14\cos4x.
\end{align*}
因此其值域为 $[1/2,1]$。下界在 $x=\pi/4+k\pi/2$ 时取得，上界在 $x=k\pi/2$ 时取得。
\end{example}

\subsection{半角公式与符号选择}

把式 \eqref{eq:trig-power-reduction} 中的 $x$ 换为 $x/2$，得到\keyterm{半角公式}（half-angle formulas）
\begin{equation}\label{eq:trig-half-square}
  \sin^2\frac x2=\frac{1-\cos x}{2},\qquad
  \cos^2\frac x2=\frac{1+\cos x}{2}.
\end{equation}
开方后写作
\begin{equation}\label{eq:trig-half-root}
  \sin\frac x2=\pm\sqrt{\frac{1-\cos x}{2}},\qquad
  \cos\frac x2=\pm\sqrt{\frac{1+\cos x}{2}}.
\end{equation}
两个符号分别由 $x/2$ 所在象限决定，不能任意选取，也不必相同。例如 $x=3\pi$ 时，$\sin(x/2)=-1$，必须取负号。

由倍角公式和式 \eqref{eq:trig-half-square} 还可得到
\begin{align}
  \tan\frac x2&=\frac{\sin x}{1+\cos x}
  &&\bigl(x\ne(2k+1)\pi\bigr),\label{eq:trig-half-tan-a}\\
  \tan\frac x2&=\frac{1-\cos x}{\sin x}
  &&\bigl(x\ne k\pi\bigr).\label{eq:trig-half-tan-b}
\end{align}
第一式的右侧为 $2\sin(x/2)\cos(x/2)/(2\cos^2(x/2))$；第二式的推导则还约去了 $\sin(x/2)$。所以二者的书写形式具有不同的适用范围：例如 $x=0$ 时第一式可用，第二式不可用。

\begin{example}[用半角求值]
因为 $\pi/8$ 是锐角，
\[
  \sin\frac\pi8=\frac{\sqrt{2-\sqrt2}}2,\qquad
  \cos\frac\pi8=\frac{\sqrt{2+\sqrt2}}2,\qquad
  \tan\frac\pi8=\frac{\sin(\pi/4)}{1+\cos(\pi/4)}=\sqrt2-1.
\]
先确定半角的符号，再开平方，可以避免把平方关系误当成函数值关系。
\end{example}

\subsection{三倍角公式}

把 $3x$ 写成 $2x+x$，展开并消去平方项，得到
\begin{equation}\label{eq:trig-triple}
  \boxed{\sin3x=3\sin x-4\sin^3x,\qquad
  \cos3x=4\cos^3x-3\cos x}.
\end{equation}
具体地，正弦展开为
\[
  2\sin x\cos^2x+(1-2\sin^2x)\sin x
  =3\sin x-4\sin^3x;
\]
余弦展开为
\[
  (2\cos^2x-1)\cos x-2\sin^2x\cos x
  =4\cos^3x-3\cos x.
\]
若 $\cos x\ne0$ 且 $1-3\tan^2x\ne0$，再取商可得
\[
  \tan3x=\frac{3\tan x-\tan^3x}{1-3\tan^2x}.
\]
三倍角公式适合处理含 $\sin^3x$、$\cos^3x$ 的表达式，或含 $x$ 与 $3x$ 的方程。

\section{积化和差公式}
\label{sec:trig-product-to-sum}

\subsection{四条公式及推导}

\keyterm{积化和差公式}（product-to-sum identities）把两个三角函数的乘积改写为和角、差角的函数值之和或差。对任意实数 $\alpha,\beta$，有
\begin{align}
  \boxed{\sin\alpha\sin\beta
  =\frac12[\cos(\alpha-\beta)-\cos(\alpha+\beta)]},\label{eq:trig-prod-ss}\\
  \boxed{\cos\alpha\cos\beta
  =\frac12[\cos(\alpha-\beta)+\cos(\alpha+\beta)]},\label{eq:trig-prod-cc}\\
  \boxed{\sin\alpha\cos\beta
  =\frac12[\sin(\alpha+\beta)+\sin(\alpha-\beta)]},\label{eq:trig-prod-sc}\\
  \boxed{\cos\alpha\sin\beta
  =\frac12[\sin(\alpha+\beta)-\sin(\alpha-\beta)]}.\label{eq:trig-prod-cs}
\end{align}

\begin{proof}
余弦的差角公式减去和角公式，右侧为 $2\sin\alpha\sin\beta$；两式相加，右侧为 $2\cos\alpha\cos\beta$。分别除以 $2$，得到前两条公式。
正弦的和角公式与差角公式相加，右侧为 $2\sin\alpha\cos\beta$；和角公式减去差角公式，右侧为 $2\cos\alpha\sin\beta$。分别除以 $2$，得到后两条公式。
\end{proof}

记忆时注意：两个正弦相乘得到余弦之差，两个余弦相乘得到余弦之和；正弦与余弦相乘得到正弦。后两式可以借助 $\sin(\beta-\alpha)=-\sin(\alpha-\beta)$ 相互转换，交换因子时须同时核对差角的顺序。

\subsection{求值、化简与求和}

\begin{example}[乘积求值]
由于 $\pi/12+5\pi/12=\pi/2$，$\pi/12-5\pi/12=-\pi/3$，
\[
  \sin\frac\pi{12}\sin\frac{5\pi}{12}
  =\frac12\left(\cos\left(-\frac\pi3\right)-\cos\frac\pi2\right)
  =\frac14.
\]
分别计算两个正弦会引入根式，先积化和差则只需计算基本特殊角。
\end{example}

\begin{example}[消去中间频率]
对任意实数 $x$，
\begin{align*}
  \sin3x\cos2x-\cos3x\sin2x
  &=\frac12(\sin5x+\sin x)-\frac12(\sin5x-\sin x)\\
  &=\sin x.
\end{align*}
这里也可以直接识别为正弦差角公式。能够看出完整的和差角结构时，直接合并更简洁；只有部分乘积时，积化和差更便于整理。
\end{example}

\begin{example}[乘积形式的有限和]
设 $n$ 为正整数，且 $\sin x\ne0$。由式 \eqref{eq:trig-prod-ss}，
\[
  2\sin x\sin(2jx)=\cos((2j-1)x)-\cos((2j+1)x).
\]
对 $j=1,\ldots,n$ 求和，中间的余弦项相互抵消，得到
\[
  \sum_{j=1}^{n}\sin(2jx)
  =\frac{\cos x-\cos((2n+1)x)}{2\sin x}.
\]
若 $\sin x=0$，原和的每一项都为零，应单独得出和为 $0$，不能使用分母为零的商式。
\end{example}

\section{和差化积公式}
\label{sec:trig-sum-to-product}

\subsection{由半和、半差统一推导}

\keyterm{和差化积公式}（sum-to-product identities）把两个三角函数值的和或差改写为乘积。对任意实数 $\alpha,\beta$，有
\begin{align}
  \boxed{\sin\alpha+\sin\beta
  =2\sin\frac{\alpha+\beta}{2}\cos\frac{\alpha-\beta}{2}},\label{eq:trig-sum-ss}\\
  \boxed{\sin\alpha-\sin\beta
  =2\cos\frac{\alpha+\beta}{2}\sin\frac{\alpha-\beta}{2}},\label{eq:trig-diff-ss}\\
  \boxed{\cos\alpha+\cos\beta
  =2\cos\frac{\alpha+\beta}{2}\cos\frac{\alpha-\beta}{2}},\label{eq:trig-sum-cc}\\
  \boxed{\cos\alpha-\cos\beta
  =-2\sin\frac{\alpha+\beta}{2}\sin\frac{\alpha-\beta}{2}}.\label{eq:trig-diff-cc}
\end{align}

\begin{proof}
令 $u=(\alpha+\beta)/2$、$v=(\alpha-\beta)/2$，则 $\alpha=u+v$、$\beta=u-v$。展开 $\sin(u+v)\pm\sin(u-v)$，交叉项抵消后分别得到 $2\sin u\cos v$ 与 $2\cos u\sin v$。展开 $\cos(u+v)\pm\cos(u-v)$，分别得到 $2\cos u\cos v$ 与 $-2\sin u\sin v$。代回 $u,v$ 即得四式。
\end{proof}

四条公式都使用角的\keyterm{半和}与\keyterm{半差}。余弦之差带负号；也可以把最后一项改写为 $2\sin((\alpha+\beta)/2)\sin((\beta-\alpha)/2)$，但不可只交换差角而保留原符号。

\subsection{求值与分式化简}

\begin{example}[和式求值]
\[
  \sin\frac\pi{12}+\sin\frac{5\pi}{12}
  =2\sin\frac\pi4\cos\left(-\frac\pi6\right)
  =\frac{\sqrt6}{2}.
\]
和差化积在半和或半差是特殊角时尤其有效。
\end{example}

\begin{example}[约分前保留定义域]
化简
\[
  F(x)=\frac{\sin3x+\sin x}{\cos3x+\cos x}.
\]
由和差化积，分子为 $2\sin2x\cos x$，分母为 $2\cos2x\cos x$。原式定义域要求 $\cos2x\cos x\ne0$，故在这个定义域上
\[
  F(x)=\tan2x.
\]
不能据此把定义域扩大成仅要求 $\cos2x\ne0$：例如 $x=\pi/2$ 时，$\tan2x$ 有定义而原分式没有定义。
\end{example}

\subsection{从和差化积到方程分解}

\begin{example}[完整保留各个因子产生的解]
求 $0\le x<2\pi$ 内方程 $\sin3x=\sin x$ 的所有解。移项并使用式 \eqref{eq:trig-diff-ss}，得
\[
  2\cos2x\sin x=0.
\]
所以 $\sin x=0$ 或 $\cos2x=0$，最终解为
\[
  x\in\left\{0,\frac\pi4,\frac{3\pi}4,\pi,
                    \frac{5\pi}4,\frac{7\pi}4\right\}.
\]
如果直接除以 $\sin x$，就会漏掉 $x=0,\pi$。求解方程时，因式分解后逐一讨论零因子比直接约去未知因子可靠。
\end{example}

\section{辅助角公式与最值}
\label{sec:trig-auxiliary}

对同一个角的正弦、余弦作线性组合时，可以利用\keyterm{辅助角公式}（auxiliary-angle formula）将其合并。

\begin{proposition}[辅助角公式]
设 $p,q$ 是不同时为零的实数，$\rho=\sqrt{p^2+q^2}$。选择实数 $\varphi$ 使
\[
  \cos\varphi=\frac p\rho,\qquad \sin\varphi=\frac q\rho,
\]
则对任意实数 $x$，
\begin{equation}\label{eq:trig-auxiliary}
  \boxed{p\sin x+q\cos x=\rho\sin(x+\varphi)}.
\end{equation}
\end{proposition}

\begin{proof}
由于 $(p/\rho)^2+(q/\rho)^2=1$，单位圆定义保证存在这样的角 $\varphi$。展开右侧的正弦和角公式，两个系数分别为 $\rho\cos\varphi=p$ 和 $\rho\sin\varphi=q$，即得等式。
\end{proof}

辅助角应由正弦与余弦两个条件共同确定，不能只用 $\tan\varphi=q/p$ 忽略象限。当 $x$ 遍历全体实数时，式 \eqref{eq:trig-auxiliary} 的值域为 $[-\rho,\rho]$；最大值在 $x+\varphi=\pi/2+2k\pi$ 时取得，最小值在 $x+\varphi=-\pi/2+2k\pi$ 时取得。如果 $x$ 被限制在一个区间内，还须检查这些角是否可达。若 $p=q=0$，原式恒为 $0$，无需引入辅助角。

\begin{example}[合并与区间最值]
因为 $\sqrt3\sin x+\cos x=2\sin(x+\pi/6)$，全体实数上的最大值为 $2$，最小值为 $-2$。
若限制 $0\le x\le\pi/2$，则 $x+\pi/6\in[\pi/6,2\pi/3]$。正弦在该区间先增后减，最大值为 $1$，两端值中较小者为 $1/2$，所以原式最大值为 $2$，在 $x=\pi/3$ 取得；最小值为 $1$，在 $x=0$ 取得。
\end{example}

\section{万能公式与半角正切代换}
\label{sec:trig-universal}

\subsection{公式的推导与适用范围}

\keyterm{万能公式}通常指以 $t=\tan(x/2)$ 为参数的三角函数有理表示，对应的代换称为\keyterm{半角正切代换}（tangent half-angle substitution）。它把正弦、余弦统一写成同一个变量的有理式。

\begin{theorem}[万能公式]
设 $x\ne(2k+1)\pi$，令 $t=\tan(x/2)$。则
\begin{equation}\label{eq:trig-universal}
  \boxed{\sin x=\frac{2t}{1+t^2},\qquad
  \cos x=\frac{1-t^2}{1+t^2}}.
\end{equation}
若进一步有 $t\ne\pm1$，则
\begin{equation}\label{eq:trig-universal-tan}
  \boxed{\tan x=\frac{2t}{1-t^2}}.
\end{equation}
\end{theorem}

\begin{proof}
假设保证 $\cos(x/2)\ne0$。把倍角公式的分子、分母同时除以 $\cos^2(x/2)$，有
\[
  \sin x=
  \frac{2\sin(x/2)\cos(x/2)}{\sin^2(x/2)+\cos^2(x/2)}
  =\frac{2t}{1+t^2},
\]
\[
  \cos x=
  \frac{\cos^2(x/2)-\sin^2(x/2)}{\cos^2(x/2)+\sin^2(x/2)}
  =\frac{1-t^2}{1+t^2}.
\]
当 $t\ne\pm1$ 时余弦非零，两式相除即得正切公式。
\end{proof}

对实数 $t$，$1+t^2$ 恒为正，因此正弦、余弦的这两个有理式本身没有实分母零点。但 $t$ 为有限实数的代换没有覆盖 $x=(2k+1)\pi$，这些角必须回到原问题单独检查。“万能”表示能将由 $\sin x,\cos x$ 组成的有理表达式转成 $t$ 的有理表达式，并不表示所有三角问题都适合这样做。

\subsection{单位圆的有理参数表示}

式 \eqref{eq:trig-universal} 也写成
\[
  (\cos x,\sin x)=\left(\frac{1-t^2}{1+t^2},\frac{2t}{1+t^2}\right).
\]
它遍历单位圆上除 $(-1,0)$ 之外的点。事实上，取过 $(-1,0)$ 的直线 $v=t(u+1)$，代入 $u^2+v^2=1$，得
\[
  (u+1)\bigl((1+t^2)u+t^2-1\bigr)=0.
\]
除去已知交点 $(-1,0)$，另一交点就是上式中的坐标。而 $t=\sin x/(1+\cos x)=\tan(x/2)$ 正是该直线的斜率。这说明漏掉的角对应同一个单位圆点。

\subsection{代换解题的步骤与例题}

应用时，先记录原式的定义域，再将 $x=(2k+1)\pi$ 单独代回检查；对其余角令 $t=\tan(x/2)$，转换为有理方程。清除分母时保留非零条件，求得 $t$ 后还原角并筛选区间。
其中 $\arctan t$ 表示位于 $(-\pi/2,\pi/2)$ 内、正切等于 $t$ 的唯一角，因此还原公式为
\begin{equation}\label{eq:trig-t-back}
  x=2\arctan t+2k\pi.
\end{equation}

\begin{example}[万能代换与遗漏角的检查]
求 $0\le x<2\pi$ 内方程 $\sin x+\cos x=-1$ 的所有解。
先检查未被代换覆盖的角 $x=\pi$，发现它满足原方程。其余角令 $t=\tan(x/2)$，则
\[
  \frac{2t+1-t^2}{1+t^2}=-1
  \quad\Longleftrightarrow\quad 2t+2=0.
\]
由于 $1+t^2>0$，清除分母是等价变形。得到 $t=-1$，即 $x=-\pi/2+2k\pi$，在指定区间内为 $3\pi/2$。故全部解为
\[
  x=\pi\quad\text{或}\quad x=\frac{3\pi}{2}.
\]
若只解关于 $t$ 的方程，会漏掉 $x=\pi$。本题也可用辅助角公式求解，两种方法可相互核验。
\end{example}

\section{基本三角方程与等价变形}
\label{sec:trig-equations}

\subsection{反三角函数与通解}

由于三角函数在全定义域上不是一一对应的，定义\keyterm{反三角函数}（inverse trigonometric functions）需要限制原函数的区间。对 $u\in[-1,1]$，$\arcsin u$ 是 $[-\pi/2,\pi/2]$ 内正弦等于 $u$ 的唯一角，$\arccos u$ 是 $[0,\pi]$ 内余弦等于 $u$ 的唯一角；$\arctan u$ 对所有实数 $u$ 有定义，其值域如式 \eqref{eq:trig-t-back} 前所述。

由单位圆的对称性与周期性，得到
\begin{align}
  \sin x=\sin\alpha
  &\Longleftrightarrow x=\alpha+2k\pi\ \text{或}\ x=\pi-\alpha+2k\pi,
  \label{eq:trig-equation-sin}\\
  \cos x=\cos\alpha
  &\Longleftrightarrow x=2k\pi\pm\alpha,
  \label{eq:trig-equation-cos}\\
  \tan x=\tan\alpha
  &\Longleftrightarrow x=\alpha+k\pi,
  \label{eq:trig-equation-tan}
\end{align}
最后一式要求两边正切有定义。前两式也可由和差化积后令每个因子为零推出；第三式来自同一斜率的单位圆点相同或互为对径点。

所以求 $\sin x=u$ 或 $\cos x=u$ 时，先判断 $|u|\le1$，再分别取 $\alpha=\arcsin u$ 或 $\alpha=\arccos u$ 代入通解。一个反三角函数值通常不是原方程的全部解。

\subsection{平方、约分与检验}

平方可能增根，约去可能为零的因子可能漏根，换元可能遗漏一部分定义域。这三种风险应分别通过回代、分类讨论和检查换元覆盖范围处理。

\begin{example}[平方产生的额外候选解]
求 $0\le x<2\pi$ 内方程 $\sin x=\cos x$ 的解。平方得到 $\sin^2x=\cos^2x$，即 $\cos2x=0$，候选角为
\[
  \frac\pi4,\quad\frac{3\pi}4,\quad\frac{5\pi}4,\quad\frac{7\pi}4.
\]
回代时 $3\pi/4$ 与 $7\pi/4$ 的正弦、余弦异号，须排除，所以解为 $\pi/4$、$5\pi/4$。也可以先说明 $\cos x=0$ 不满足原式，再除以 $\cos x$ 得 $\tan x=1$，直接使用式 \eqref{eq:trig-equation-tan}。
\end{example}

\section{三角形记号与正弦定理}
\label{sec:trig-sine-law}

\subsection{统一记号与面积公式}

下文只讨论欧氏平面中的非退化三角形 $ABC$。也用 $A,B,C$ 表示三个内角的大小，约定
\[
  a=|BC|,\qquad b=|CA|,\qquad c=|AB|,\qquad
  A,B,C\in(0,\pi),\quad A+B+C=\pi.
\]
边长 $a,b,c$ 均为正，并满足严格三角不等式。记面积为 $\Delta$，半周长为 $s=(a+b+c)/2$，外接圆半径为 $R$，内切圆半径为 $r$。

以 $AB$ 为底，点 $C$ 到直线 $AB$ 的距离为 $b\sin A$，无论 $A$ 为锐角还是钝角，该距离均为正。因此
\begin{equation}\label{eq:trig-area}
  \boxed{\Delta=\frac12bc\sin A
  =\frac12ca\sin B=\frac12ab\sin C}.
\end{equation}
后两式由轮换顶点得到。

\subsection{正弦定理及证明}

\begin{theorem}[正弦定理]
对上述三角形，有\keyterm{正弦定理}（law of sines）
\begin{equation}\label{eq:trig-sine-law}
  \boxed{\frac a{\sin A}=\frac b{\sin B}=\frac c{\sin C}=2R}.
\end{equation}
\end{theorem}

\begin{proof}
由面积公式，$\sin A=2\Delta/(bc)$，从而
\[
  \frac a{\sin A}=\frac{abc}{2\Delta}.
\]
轮换字母得到三个比值相等。
为确定比值与外接圆的关系，设弦 $BC$ 对应的较小圆心角为 $\theta\in(0,\pi]$。在圆心与弦端点形成的等腰三角形中作高，或在 $\theta=\pi$ 时直接使用直径，得到 $a=2R\sin(\theta/2)$。由圆周角定理，$A$ 等于 $\theta/2$ 或 $\pi-\theta/2$，两者的正弦相同。因此 $a=2R\sin A$，即公共比值等于 $2R$。
\end{proof}

由证明同时得到
\begin{equation}\label{eq:trig-circum-area}
  \Delta=\frac{abc}{4R}.
\end{equation}
正弦定理适合已知一组对边、对角关系的情形，但由一个正弦值反求三角形内角时，应注意锐角与其补角具有相同正弦。

\begin{example}[已知两角和一边]
设 $A=\pi/6$、$B=\pi/4$、$a=2$，则
\[
  C=\pi-A-B=\frac{7\pi}{12},\qquad
  2R=\frac2{\sin(\pi/6)}=4.
\]
所以
\[
  b=4\sin\frac\pi4=2\sqrt2,\qquad
  c=4\sin\frac{7\pi}{12}=\sqrt6+\sqrt2,\qquad R=2.
\]
角和先确定第三角，正弦定理再确定其余边，三角形的大小与形状均被确定。
\end{example}

\subsection{两边及一边对角：零解、一解或两解}
\label{subsec:trig-ssa}

设已知 $a,b,A$，其中 $a,b>0$、$0<A<\pi$。由正弦定理令
\[
  q=\frac{b\sin A}{a}=\sin B.
\]
若 $q>1$，无解；若 $0<q\le1$，令 $B_1=\arcsin q\in(0,\pi/2]$，候选角为 $B_1$ 与 $B_2=\pi-B_1$。只保留满足 $A+B_i<\pi$ 的候选；当 $q=1$ 时两者相同，只计一次。每个有效候选确定 $C_i=\pi-A-B_i$ 及 $c_i=a\sin C_i/\sin A$。这给出完整的判定方法。

也可按长度比较快速分类。令 $h=b\sin A$，则
\begin{itemize}
  \item 若 $A<\pi/2$：$a<h$ 时无解；$a=h$ 时一解；$h<a<b$ 时两解；$a\ge b$ 时一解。
  \item 若 $A\ge\pi/2$：$a>b$ 时一解，$a\le b$ 时无解。
\end{itemize}
这些条件可以从候选角法核验。锐角 $A$ 下，$a>h$ 保证存在锐角候选 $B_1$，且 $A+B_1<\pi$；钝角候选 $\pi-B_1$ 有效当且仅当 $B_1>A$，由正弦在锐角区间的严格递增性，这等价于 $a<b$。非锐角 $A$ 下，只有锐角候选可能有效，条件 $B_1<\pi-A$ 等价于 $a>b$；等号会导致第三角为零，不能计为三角形。

\begin{example}[出现两个三角形]
设 $A=\pi/6$、$a=1$、$b=\sqrt3$，则 $\sin B=\sqrt3/2$，从而
\[
  B=\frac\pi3\quad\text{或}\quad B=\frac{2\pi}3.
\]
两者分别给出 $C=\pi/2$ 与 $C=\pi/6$，都满足内角条件。又 $a/\sin A=2$，所以对应的第三边分别为 $c=2$ 与 $c=1$。若只取反正弦的主值，就会漏掉第二个三角形。
\end{example}

\section{余弦定理与三角形判定}
\label{sec:trig-cosine-law}

\subsection{余弦定理及坐标证明}

\begin{theorem}[余弦定理]
对非退化三角形 $ABC$，有\keyterm{余弦定理}（law of cosines）
\begin{align}
  a^2&=b^2+c^2-2bc\cos A,\label{eq:trig-cosine-law-a}\\
  b^2&=c^2+a^2-2ca\cos B,\notag\\
  c^2&=a^2+b^2-2ab\cos C.\notag
\end{align}
\end{theorem}

\begin{proof}
取 $A=(0,0)$、$B=(c,0)$，使 $C$ 位于横轴上方，则 $C=(b\cos A,b\sin A)$。由两点间距离公式，
\[
  a^2=(b\cos A-c)^2+b^2\sin^2A
  =b^2+c^2-2bc\cos A.
\]
这个坐标表示对锐角、直角、钝角均成立。轮换顶点得到其余两式。
\end{proof}

当 $A=\pi/2$ 时，余弦项为零，退化为勾股定理。已知两边及夹角时，用式 \eqref{eq:trig-cosine-law-a} 求对边；已知三边时，用
\begin{equation}\label{eq:trig-cosine-angle}
  \boxed{\cos A=\frac{b^2+c^2-a^2}{2bc}}
\end{equation}
求角。由于余弦在 $(0,\pi)$ 上严格递减，三角形内角由余弦值唯一确定，不存在正弦反求角时的补角歧义。

\begin{example}[两边及夹角求第三边]
设 $b=3$、$c=5$、$A=\pi/3$，则
\[
  a^2=9+25-2\cdot3\cdot5\cdot\frac12=19,
  \qquad a=\sqrt{19}.
\]
面积为 $\Delta=bc\sin A/2=15\sqrt3/4$。若还需 $B$，可用
\[
  B=\arccos\frac{a^2+c^2-b^2}{2ac}
   =\arccos\frac7{2\sqrt{19}},\qquad C=\pi-A-B.
\]
\end{example}

\subsection{由三边判断角的类型}

式 \eqref{eq:trig-cosine-angle} 的分母为正，所以
\[
\begin{array}{lll}
  a^2<b^2+c^2 &\Longleftrightarrow \cos A>0 &\Longleftrightarrow A<\pi/2,\\
  a^2=b^2+c^2 &\Longleftrightarrow \cos A=0 &\Longleftrightarrow A=\pi/2,\\
  a^2>b^2+c^2 &\Longleftrightarrow \cos A<0 &\Longleftrightarrow A>\pi/2.
\end{array}
\]
判断整个三角形是锐角、直角还是钝角三角形时，应让 $a$ 表示最长边，并先检查三角不等式。例如边长 $2,3,4$ 能组成三角形，且 $4^2>2^2+3^2$，所以是钝角三角形。

这里使用了大边对大角的性质，也可用正弦定理与和差化积证明：
\[
  a-b=2R(\sin A-\sin B)
  =4R\cos\frac{A+B}{2}\sin\frac{A-B}{2}.
\]
由于 $0<(A+B)/2<\pi/2$，余弦因子为正；且 $(A-B)/2\in(-\pi/2,\pi/2)$，其正弦与 $A-B$ 同号。因此 $a-b$ 与 $A-B$ 同号，即 $a>b$ 当且仅当 $A>B$。

\section{面积、半角与其他三角形公式}
\label{sec:trig-triangle-more}

\subsection{海伦公式及其推导}

\begin{theorem}[海伦公式]
设 $s=(a+b+c)/2$，则三角形面积满足\keyterm{海伦公式}（Heron's formula）
\begin{equation}\label{eq:trig-heron}
  \boxed{\Delta=\sqrt{s(s-a)(s-b)(s-c)}}.
\end{equation}
\end{theorem}

\begin{proof}
由面积公式与余弦定理，
\begin{align*}
  16\Delta^2
  &=4b^2c^2\sin^2 A\\
  &=4b^2c^2-(b^2+c^2-a^2)^2\\
  &=\bigl(a^2-(b-c)^2\bigr)\bigl((b+c)^2-a^2\bigr)\\
  &=(a-b+c)(a+b-c)(b+c-a)(a+b+c)\\
  &=16s(s-a)(s-b)(s-c).
\end{align*}
严格三角不等式保证根号内各因子为正，且 $\Delta>0$，故取正平方根。
\end{proof}

把内心与三个顶点连接，三个小三角形到各边的高都是 $r$，面积相加得到
\begin{equation}\label{eq:trig-inradius}
  \Delta=\frac12r(a+b+c)=rs.
\end{equation}
结合式 \eqref{eq:trig-circum-area}，已知三边时可依次求出 $\Delta$、$r=\Delta/s$ 与 $R=abc/(4\Delta)$。

\begin{example}[由三边求面积与两种半径]
设三边为 $13,14,15$，则 $s=21$，由海伦公式
\[
  \Delta=\sqrt{21\cdot8\cdot7\cdot6}=84.
\]
所以 $r=84/21=4$，$R=13\cdot14\cdot15/(4\cdot84)=65/8$。这些值分别描述三角形的面积、内切圆大小与外接圆大小。
\end{example}

\subsection{三角形中的半角公式}

由于三角形内角 $A\in(0,\pi)$，其半角 $A/2$ 一定为锐角，半角公式均取正号。代入式 \eqref{eq:trig-cosine-angle}，有
\[
  \sin^2\frac A2
  =\frac{2bc-b^2-c^2+a^2}{4bc}
  =\frac{(s-b)(s-c)}{bc},
\]
\[
  \cos^2\frac A2
  =\frac{2bc+b^2+c^2-a^2}{4bc}
  =\frac{s(s-a)}{bc}.
\]
因此
\begin{equation}\label{eq:trig-triangle-half}
  \boxed{\sin\frac A2=\sqrt{\frac{(s-b)(s-c)}{bc}},\qquad
  \cos\frac A2=\sqrt{\frac{s(s-a)}{bc}}}.
\end{equation}
两式相除，再用海伦公式及 $\Delta=rs$，得到
\begin{equation}\label{eq:trig-triangle-half-tan}
  \tan\frac A2
  =\sqrt{\frac{(s-b)(s-c)}{s(s-a)}}=\frac r{s-a}.
\end{equation}
关于 $B/2$、$C/2$ 的公式由同时轮换边角字母得到；所有分母均因非退化条件而为正。

\subsection{内角和带来的恒等式}

三角形内角并非三个独立变量。由 $A+B=\pi-C$ 与和差化积，
\begin{align*}
  \sin A+\sin B+\sin C
  &=2\cos\frac C2\left(\cos\frac{A-B}{2}+\sin\frac C2\right)\\
  &=2\cos\frac C2\left(\cos\frac{A-B}{2}+\cos\frac{A+B}{2}\right)\\
  &=4\cos\frac A2\cos\frac B2\cos\frac C2.
\end{align*}
这是和差化积在有角和约束时的典型应用。

若 $A,B,C$ 均不为直角，则还有
\begin{equation}\label{eq:trig-triangle-tan}
  \tan A+\tan B+\tan C=\tan A\tan B\tan C.
\end{equation}
证明时展开 $\sin(A+B+C)=\sin\pi=0$，得到
\[
  \sin A\cos B\cos C+\cos A\sin B\cos C
  +\cos A\cos B\sin C-\sin A\sin B\sin C=0,
\]
再除以非零的 $\cos A\cos B\cos C$ 即得。若三角形含直角，式 \eqref{eq:trig-triangle-tan} 中有正切未定义，不能使用。

\section{选用公式的策略与综合练习}

\subsection{按表达式结构选择变形}

公式之间的主要推导关系是：单位圆与旋转给出和差角公式；令两角相等得到倍角公式；解出平方项得到降幂与半角公式；对和差角公式相加、相减得到积化和差；改用半和与半差变量得到和差化积；以半角正切为参数得到万能公式。

实际解题时可按以下结构选择工具：
\begin{itemize}
  \item 遇到两个三角函数相乘，希望求和或消去某些角时，考虑积化和差；遇到两个函数值相加、相减，希望求零点或约分时，考虑和差化积。
  \item 遇到平方、四次幂时，先用平方关系与降幂；遇到 $x$、$2x$ 或 $3x$ 同时出现时，尝试统一到同一个角。
  \item 遇到 $p\sin x+q\cos x$ 的求值、方程或最值问题，先考虑辅助角；遇到正弦、余弦交织的有理式，再考虑万能代换。
  \item 解三角形时，已知两角和一边用角和及正弦定理；已知两边及夹角用余弦定理；已知三边先检查三角不等式，再用余弦定理；已知两边及一边对角按第 \ref{subsec:trig-ssa} 节检查解的个数。
\end{itemize}
每次变形都应记录定义域，反求角时保留所有分支，最终按题设区间筛选。公式给出的表达式相同，并不自动表示两侧作为函数具有相同定义域。

\subsection{练习}

\begin{enumerate}
  \item 将 $\cos5x\cos x$ 化为和式，并将 $\cos5x-\cos x$ 化为积式。
  \item 求 $\sin(\pi/12)\cos(\pi/12)$ 的精确值。
  \item 求 $0\le x<2\pi$ 内方程 $\cos2x+\cos x=0$ 的所有解。
  \item 求 $3\sin x-4\cos x$ 在全体实数上的最大值和最小值，并写出最大值的取等条件。
  \item 用万能代换求 $0\le x<2\pi$ 内方程 $2\sin x=1+\cos x$ 的所有解。
  \item 三角形中 $a=7$、$b=8$、$c=9$，求 $\cos A$、面积 $\Delta$ 及外接圆半径 $R$。
  \item 已知 $A=\pi/6$、$b=2$，分别讨论 $a=1/2,1,3/2,2$ 时三角形解的个数。
\end{enumerate}

\subsection{参考解答}

\begin{enumerate}
  \item 由积化和差与和差化积，分别为
  \[
    \cos5x\cos x=\frac12(\cos6x+\cos4x),\qquad
    \cos5x-\cos x=-2\sin3x\sin2x.
  \]
  \item 由倍角公式，原式为 $\tfrac12\sin(\pi/6)=1/4$。
  \item 和差化积得 $2\cos(3x/2)\cos(x/2)=0$。第一因子为零给出 $x=\pi/3+2k\pi/3$，第二因子为零给出 $x=\pi+2k\pi$。筛选并去重后，解为 $\pi/3,\pi,5\pi/3$。
  \item 取 $\varphi\in(-\pi/2,0)$ 满足 $\cos\varphi=3/5$、$\sin\varphi=-4/5$，则原式为 $5\sin(x+\varphi)$。最大值为 $5$，最小值为 $-5$；最大值在 $x=\pi/2-\varphi+2k\pi$ 时取得。
  \item 先检查 $x=\pi$，它满足方程。对其余角作代换，得到 $4t/(1+t^2)=2/(1+t^2)$，故 $t=1/2$。全部解为 $x=2\arctan(1/2)$ 与 $x=\pi$。
  \item 三角不等式成立。余弦定理给出 $\cos A=(64+81-49)/(2\cdot8\cdot9)=2/3$。半周长 $s=12$，海伦公式给出 $\Delta=\sqrt{12\cdot5\cdot4\cdot3}=12\sqrt5$，因此 $R=7\cdot8\cdot9/(48\sqrt5)=21\sqrt5/10$。
  \item 此时 $h=b\sin A=1$，且 $A$ 为锐角。四种情形依次为 $a<h$、$a=h$、$h<a<b$、$a=b$，解的个数依次为 $0,1,2,1$。
\end{enumerate}

\end{document}
